\chapter{Number theory}

\section{模运算}
	\kactlimport{ModularArithmetic.h}

	\begin{lecture}{模逆元}
		\[
		\begin{array}{c}
		\begin{array}{c}
		\dfrac{a}{b} \equiv x \pmod p\\
		\text{假设 } \exists b_{\mathrm{inv}},\ \mathrm{s.t.}\ a \cdot b_{\mathrm{inv}} \equiv x \pmod p
		\end{array}\\[1em]
		\begin{array}{rl}
		\Rightarrow & a \equiv bx \pmod p\\
		\Rightarrow & a \cdot b_{\mathrm{inv}} \cdot b \equiv bx \pmod p\\
		\Rightarrow & a \cdot (b_{\mathrm{inv}} \cdot b - 1) \equiv 0 \pmod p\\
		\Rightarrow & b_{\mathrm{inv}} \cdot b \equiv 1 \pmod p
		\end{array}\\[1em]
		\mathrm{inv}\ \text{存在} \Leftrightarrow \gcd(b,p) = 1\\[1em]
		\begin{array}{rl}
		b \cdot b_{\mathrm{inv}} = k \cdot p + 1,\ 
		g = \gcd(b,p) \Rightarrow
		& g \cdot b_1 \cdot b_{\mathrm{inv}} = k \cdot g \cdot p_1 + 1\\
		\Rightarrow
		& g \cdot (b_1 \cdot b_{\mathrm{inv}} - k \cdot p_1) = 1\\
		\text{数域为 } \mathbb{Z} \Rightarrow
		& g = 1,\ b_1 \cdot b_{\mathrm{inv}} - k \cdot p_1 = 1
		\end{array}
		\end{array}
		\]
	\end{lecture}
	\begin{lecture}{2}
		\[
\begin{aligned}
&\frac{a}{b} \equiv x \pmod p, \qquad
\text{假设 } \exists\, b_{\mathrm{inv}},\ 
a b_{\mathrm{inv}} \equiv x \pmod p .
\\[0.8em]
&\begin{aligned}
a &\equiv bx \pmod p \\
a b_{\mathrm{inv}} b &\equiv bx \pmod p \\
a\bigl(b_{\mathrm{inv}}b - 1\bigr) &\equiv 0 \pmod p \\
b_{\mathrm{inv}}b &\equiv 1 \pmod p
\end{aligned}
\\[1em]
&\therefore\quad
\mathrm{inv}\ \text{存在}
\Longleftrightarrow
\gcd(b,p)=1 .
\\[1em]
&\begin{aligned}
b b_{\mathrm{inv}} &= kp + 1, 
\qquad g=\gcd(b,p),                                    \\[0.3em]
g b_1 b_{\mathrm{inv}} &= k g p_1 + 1,                  \\
g\bigl(b_1 b_{\mathrm{inv}} - k p_1\bigr) &= 1 .
\end{aligned}
\\[1em]
&\text{数域为 } \mathbb{Z}
\quad\Longrightarrow\quad
g=1,\qquad
b_1 b_{\mathrm{inv}} - k p_1 = 1 .
\end{aligned}
\]
	\end{lecture}
	\begin{lecture}{3}
		\[
\begin{aligned}
&\frac{a}{b} \equiv x \pmod p,
\qquad
\text{假设 } \exists\, b_{\mathrm{inv}},\
a b_{\mathrm{inv}} \equiv x \pmod p .
\\[0.8em]
\Rightarrow\quad
&a \equiv bx \pmod p \\
\Rightarrow\quad
&a b_{\mathrm{inv}} b \equiv bx \pmod p \\
\Rightarrow\quad
&a\bigl(b_{\mathrm{inv}}b - 1\bigr) \equiv 0 \pmod p \\
\Rightarrow\quad
&b_{\mathrm{inv}}b \equiv 1 \pmod p
\\[1em]
&\mathrm{inv}\ \text{存在}
\Longleftrightarrow
\gcd(b,p)=1 .
\\[1em]
&b b_{\mathrm{inv}} = kp+1,\qquad g=\gcd(b,p)
\\
\Rightarrow\quad
&g b_1 b_{\mathrm{inv}} = k g p_1 + 1 \\
\Rightarrow\quad
&g\bigl(b_1 b_{\mathrm{inv}} - k p_1\bigr) = 1 \\
\text{数域为 } \mathbb{Z}
\Rightarrow\quad
&g=1,\qquad b_1 b_{\mathrm{inv}} - k p_1 = 1 .
\end{aligned}
\]
	\end{lecture}

	\kactlimport{ModInverse.h}
	\kactlimport{ModPow.h}
	\kactlimport{ModLog.h}
	\kactlimport{ModSum.h}
	\kactlimport{ModMulLL.h}
	\kactlimport{ModSqrt.h}

\section{Primality}
	\kactlimport{FastEratosthenes.h}
	\kactlimport{MillerRabin.h}
	\kactlimport{Factor.h}

\section{Divisibility}
	\kactlimport{euclid.h}
	% \kactlimport{Euclid.java}
	\kactlimport{CRT.h}

	\subsection{Bézout's identity}
	For $a \neq $, $b \neq 0$, then $d=gcd(a,b)$ is the smallest positive integer for which there are integer solutions to
	$$ax+by=d$$
	If $(x,y)$ is one solution, then all solutions are given by
	$$\left(x+\frac{kb}{\gcd(a,b)}, y-\frac{ka}{\gcd(a,b)}\right), \quad k\in\mathbb{Z}$$

	\kactlimport{phiFunction.h}

\section{Fractions}
	\kactlimport{ContinuedFractions.h}
	\kactlimport{FracBinarySearch.h}

\section{Pythagorean Triples}
 The Pythagorean triples are uniquely generated by
 \[ a=k\cdot (m^{2}-n^{2}),\ \,b=k\cdot (2mn),\ \,c=k\cdot (m^{2}+n^{2}), \]
 with $m > n > 0$, $k > 0$, $m \bot n$, and either $m$ or $n$ even.

\section{Primes}
	$p=962592769$ is such that $2^{21} \mid p-1$, which may be useful. For hashing
	use 970592641 (31-bit number), 31443539979727 (45-bit), 3006703054056749
	(52-bit). There are 78498 primes less than 1\,000\,000.

	Primitive roots exist modulo any prime power $p^a$, except for $p = 2, a > 2$, and there are $\phi(\phi(p^a))$ many.
	For $p = 2, a > 2$, the group $\mathbb Z_{2^a}^\times$ is instead isomorphic to $\mathbb Z_2 \times \mathbb Z_{2^{a-2}}$.

\section{Estimates}
	$\sum_{d|n} d = O(n \log \log n)$.

	The number of divisors of $n$ is at most around 100 for $n < 5e4$, 500 for $n < 1e7$, 2000 for $n < 1e10$, 200\,000 for $n < 1e19$.

\section{Mobius Function}
\[
	\mu(n) = \begin{cases} 0 & n \textrm{ is not square free}\\ 1 & n \textrm{ has even number of prime factors}\\ -1 & n \textrm{ has odd number of prime factors}\\\end{cases}
\]
  Mobius Inversion:
  \[ g(n) = \sum_{d|n} f(d) \Leftrightarrow f(n) = \sum_{d|n} \mu(d)g(n/d) \]
  Other useful formulas/forms:

  $ \sum_{d | n} \mu(d) = [ n = 1] $ (very useful)

  $ g(n) = \sum_{n|d} f(d) \Leftrightarrow f(n) = \sum_{n|d} \mu(d/n)g(d)$

 $ g(n) = \sum_{1 \leq m \leq n} f(\left\lfloor\frac{n}{m}\right \rfloor ) \Leftrightarrow f(n) = \sum_{1\leq m\leq n} \mu(m)g(\left\lfloor\frac{n}{m}\right\rfloor)$
